\addcontentsline{toc}{section}{СПИСОК ИСПОЛЬЗОВАННЫХ ИСТОЧНИКОВ}

\begin{thebibliography}{50}
\bibitem{ISOP} Каракчиев Д.А., Шевченко К.Н., Томакова Р.А. Интеллектуальная система оптимизации производственной деятельности на среднем и крупном предприятии на основе автоматизированных информационных систем // Известия Юго-Западного государственного университета. Серия: Управление, вычислительная техника, информатика. Медицинское приборостроение. – 2026. – Т. 16, № 1. – С. 50–63. –  URL: https://doi.org/10.21869/2223-1536-2026-16-1-50-63 – Текст : электронный.
\bibitem{machinostroy} Загидуллин Р. Р. Управление машиностроительным производством с помощью систем MES, APS, ERP. Полная версия : Монография / Загидуллин Р. Р. - Старый Оскол : ТНТ, 2021. – 416 с. - ISBN 978-5-94178-282-6. Текст : электронный // ЭБС ТНТ [сайт]. – URL: https://www.tnt-ebook.ru/library/book/548 (дата обращения: 02.06.2026).
\bibitem{borovkov} Боровков А. И. Интеллектуальные системы управления : учебное пособие / А. И. Боровков. – Москва : ИНФРА-М, 2022. – 345 с. – ISBN 978-5-16-017006-9. – Текст : непосредственный.  
\bibitem{neuralnetworks} Хайкин  С. Нейронные сети: полный курс : учебное пособие / С. Хайкин. – Москва : Вильямс, 2018. – 1104 с. – ISBN 978-5-8459-2101-0. – Текст : непосредственный.
\bibitem{python} Лутц  М. Изучаем Python : учебное пособие  / М. Лутц. –  5-е издание – Санкт-Петербург : Питер, 2019. – 1584 с. – ISBN 978-5-4461-0705-9. – Текст : непосредственный.
\bibitem{deeplearning} Гудфеллоу И., Бенджио Ю., Курвилль А. Глубокое обучение : учебное пособие / И. Гудфеллоу, Ю. Бенджио. –  Москва : ДМК Пресс, 2017. – 652 с. – ISBN 978-5-97060-487-9. – Текст : непосредственный.    
\bibitem{machinelearning} Мерфи К. Машинное обучение: вероятностный подход : учебное пособие / К. Мерфи. – Москва : ДМК Пресс, 2018. – 704 с. – ISBN 978-5-97060-212-7. – Текст: непосредственный.
\bibitem{pythonmachinelearning} Рашка С., Мирджалили, В. Python и машинное обучение : практическое пособие / С. Рашка, В. Мирджалили. – Москва : ДМК Пресс, 2018. – 418 с. – ISBN 978-5-97060-310-0. – Текст : непосредственный.
\bibitem{predstavlenieznaniy} Сосинская С. С. Представление знаний в информационной системе. Методы искусственного интеллекта и представления знаний : Учебное пособие / Сосинская С. С. – Старый Оскол : ТНТ, 2022. – 216 с. - ISBN 978-5-94178-254-3. Текст : электронный // ЭБС ТНТ [сайт]. – URL: https://www.tnt-ebook.ru/library/book/497 (дата обращения: 23.05.2026).
\bibitem{keras} Чоллет Ф. Глубокое обучение на Python : учебное пособие / Ф. Чоллет. – Москва : ДМК Пресс, 2018. – 304 с. – ISBN 978-5-97060-409-1. – Текст : непосредственный.
\bibitem{nechetkayainf} Еременко Ю. И. Регулирование, управление и принятие решений в условиях нечёткой информации : Учебное пособие / Еременко Ю. И. – Старый Оскол : ТНТ, 2020. – 256 с. - ISBN 978-5-94178-614-5. Текст : электронный // ЭБС ТНТ [сайт]. – URL: https://www.tnt-ebook.ru/library/book/530 (дата обращения: 12.05.2026). 
\bibitem{advancedpython} Бейдер Д. Python Tricks: A Buffet of Awesome Python Features : учебное пособие / Д. Бейдер. – Москва : ДМК Пресс, 2021. – 300 с. – ISBN 978-5-97060-999-7. – Текст : непосредственный.
\bibitem{practicalml} Герон О. Практическое машинное обучение с Scikit-Learn и TensorFlow : учебное пособие / О. Герон. – Москва : ДМК Пресс, 2019. – 572 с. – ISBN 978-5-97060-524-1. – Текст : непосредственный.
\bibitem{neuralnetspython} Нильсен М. Нейронные сети и глубокое обучение : учебное пособие / М. Нильсен. – Москва : ДМК Пресс, 2021. – 250 с. – ISBN 978-5-97060-777-1. – Текст : непосредственный.
\bibitem{uml2} Джеймс Р. UML 2.0. Объектно-ориентированное моделирование и разработка : практическое пособие / Р. Джеймс, Б. Майкл. – 2-е изд. – Санкт-Петербург : Питер, 2021. – 542 с. – ISBN 978-5-4461-9428-5. – Текст : непосредственный.
\bibitem{oopanalyz} Зайцев, М. Г. Объектно-ориентированный анализ и программирование : учебное пособие / М. Г. Зайцев. – Новосибирск : изд-во НГТУ, 2017. – 84 с. – ISBN 978-5-04-112962-0. – Текст : непосредственный.
\bibitem{interface} Мандел, Т. Разработка пользовательского интерфейса : учебное пособие / Т. Мандел. – Москва : ДМК Пресс, 2019. – 420 с. – ISBN 978-5-04-195060-6. – Текст : непосредственный.
\bibitem{setiavtomat} Сети автоматизации : Учебное пособие / Лыков А. Н., Катаев Р. В., Бочкарев С. В., Петроченков А. Б. – Старый Оскол : ТНТ, 2020. – 432 с. - ISBN 978-5-94178-674-9. Текст : непосредственный.
\bibitem{strategia} Свиридова С. В. Разработка теоретического и научно-методического обеспечения стратегического инновационного развития промышленных предприятий : дис. ... д-ра экон. наук : 08.00.05 / С. В. Свиридова ; науч. конс. д-р экон. наук, проф. Ю. П. Анисимова ; Воронеж. гос. техн. ун-т. – Воронеж, 2016. – 420 с. – Библиогр.: с. 355–383. – Текст : непосредственный.
\bibitem{kachestvo} Петровский Э. А. Управление качеством производственных и технологических систем : Учебник / Петровский Э. А. – 2-е издание, переработанное и дополненное – Старый Оскол : ТНТ, 2026. – 351 с. - ISBN 978-5-94178-547-6. Текст : электронный // ЭБС ТНТ [сайт]. – URL: https://www.tnt-ebook.ru/library/book/146 (дата обращения: 01.03.2026).
\bibitem{teoriaNS} Томакова Р.А.  Основы теории нейрокомпьютерных систем : учебное пособие / Р. А. Томакова ; Юго-Зап. гос. ун-т. - Курск : [б. и.], 2021. - 135 с. - ISBN 978-5-907555-36-5 : Б. ц. - Текст : непосредственный.
\bibitem{methodbase} Томакова Р.А.  Методологические основы научных исследований : учебное пособие / Р. А. Томакова, В. И. Томаков ; Юго-Зап. гос. ун-т. - Курск : ЮЗГУ, 2017. - 204 с. - Библиогр.: с. 199-203. - ISBN 978-5-7681-1210-3. - Текст : непосредственный.
\bibitem{statanaliz} Программные системы статистического анализа: обнаружение закономерностей в данных с использованием системы R и языка Python : учебное пособие / В. М. Волкова, М. А. Семенова, Е. С. Четвертакова, С. С. Вожов. - Новосибирск : Новосибирский государственный технический университет, 2017. - 74 с. : ил., табл. - URL: https://biblioclub.ru/index.php?page=book\&id=576496 (дата обращения: 02.03.2022) . - Режим доступа: по подписке. - Библиогр.: с. 48. - ISBN 978-5-7782-3183-2 : Б. ц. - Текст : электронный.
\bibitem{pyyavu} Шелудько, В. М.  Язык программирования высокого уровня Python: функции, структуры данных, дополнительные модули : учебное пособие / В. М. Шелудько ; Министерство науки и высшего образования РФ ; Федеральное государственное автономное образовательное учреждение высшего образования «Южный федеральный университет» ; Институт компьютерных технологий и информационной безопасности. - Ростов-на-Дону|Таганрог : Издательство Южного федерального университета, 2017. - 108 с. : ил. - URL: http://biblioclub.ru/index.php?page=book\&id=500060 (дата обращения: 24.04.2026) . - Режим доступа: по подписке. - Библиогр. в кн. - ISBN 978-5-9275-2648-2 : Б. ц. - Текст : электронный.
\bibitem{prinreshen} Демидова Л.А.  Принятие решений в условиях неопределенности : монография / Л. А. Демидова. - 2-е изд., перераб. - Москва : Телеком, 2016. - 289 с. - ISBN 978-5-9912-0513-9 - Текст : непосредственный.
\bibitem{teorians2} Яхъяева Г. Э. Основы теории нейронных сетей : учебное пособие / Г. Э. Яхъяева. – 2-е изд., испр. – Москва : Национальный Открытый Университет «ИНТУИТ», 2016. – 200 с. – (Основы информационных технологий). – URL: http://biblioclub.ru/index.php?page=book\&id=429110 (дата обращения: 01.06.2026). – ISBN 978-5-94774-818-5 : Б. ц. – Текст : электронный.
\bibitem{naukaodata} Келлехер, Д.    Наука о данных: базовый курс : учебное пособие / Д. Келлехер, Б. Тирни ; науч. ред. З. Мамедьяров ; пер. с англ. М. Белоголовский. - Москва : Альпина Паблишер, 2020. - 224 с. - URL: http://biblioclub.ru/index.php?page=book\&id=598235 (дата обращения: 09.02.2026) . - Режим доступа: по подписке. - ISBN 978-5-9614-3170-4 : Б. ц. - Текст : электронный.
\bibitem{matlab} Кассим К.Д.А.   Моделирование систем искусственного интеллекта в среде Matlab и Fuzzytech : учебное пособие / К. Д. А. Кассим, С. А. Филист, О. В. Шаталова. - Курск : Деловая полиграфия, 2016. - 186 с. - Библиогр.: с. 185. - ISBN 978-5-9908582-8-2. - Текст : непосредственный.
\bibitem{MLP} Д. С. Пономарев Нейросетевые методы прогнозирования в производственном секторе: MLP или RNN? // Информатика. Экономика. Управление / Informatics. Economics. Management. 2025. №2. URL: https://cyberleninka.ru/article/n/neyrosetevye-metody-prognozirovaniya-v-proizvodstvennom-sektore-mlp-ili-rnn (дата обращения: 08.06.2026).
\bibitem{allfill} Программирование, тестирование, проектирование, нейросети, технологии аппаратно‐программных средств (практические задания и способы их решения) : учебник / С. В. Веретехина, К. С. Кармицкий, Д. Д. Лукашин [и др.]. - Москва : Директ-Медиа, 2022. - 144 с. - URL: https://biblioclub.ru/index.php?page=book\&id=694782 (дата обращения: 10.01.2023) . - Режим доступа: по подписке. - Библиогр. в кн. - ISBN 978-5-4499-3321-8 : Б. ц. - Текст : электронный.
\bibitem{intsysandtech} Интеллектуальные системы и технологии : учебное пособие / С. П. Ющенко [и др.] ; Юго-Зап. гос. ун-т. - Курск : Университетская книга, 2018. - 226 с. : ил. - Библиогр.: с. 238-240 (32 назв.). - ISBN 978-5-907138-22-3 - Текст : непосредственный.
\bibitem{sidorkina} Сидоркина И.Г.  Системы искусственного интеллекта : учебное пособие / И. Г. Сидоркина. - Москва : КНОРУС, 2016. - 246 с. : рис. - Библиогр.: с. 244-245. - ISBN 978-5-406-04876-4 . - Текст : непосредственный.
\bibitem{pugacheva} Пугачева Дарья Борисовна, Юдина Майя Викторовна Исследование программных решений для определения оптимального решения по заданным параметрам // Computational nanotechnology. 2024. №5. URL: https://cyberleninka.ru/article/n/issledovanie-programmnyh-resheniy-dlya-opredeleniya-optimalnogo-resheniya-po-zadannym-parametram (дата обращения: 08.06.2026).
\bibitem{kachestvo2} Лютов А. Г. Управление качеством. Интегрированные информационно-аналитические системы : Учебник / Лютов А. Г., Огородов В. А., Чугунова О. И. 1– Старый Оскол : ТНТ, 2025. – 284 с. - ISBN 978-5-94178-868-2. Текст : непосредственный.
\bibitem{sidorkina}Лазарсон Э. В. Теория и методы решения многовариантных неформализованных задач выбора : Монография / Лазарсон Э. В. 1– Старый Оскол : ТНТ, 2022. – 240 с. - ISBN 978-5-94178-042-6. -  Текст : электронный // ЭБС ТНТ сайт. – URL: https://www.tnt-ebook.ru/library/book/476 (дата обращения: 02.06.2026).
\bibitem{isakova} Исакова, А. И.    Основы информационных технологий : учебное пособие / А. И. Исакова. - Томск : ТУСУР, 2016. - 206 с. : ил. - URL: http://biblioclub.ru/index.php?page=book\&id=480808 (дата обращения: 18.12.2025) . - Режим доступа: по подписке. - Библиогр.: с. 197-198. - Б. ц. - Текст : электронный.
\bibitem{gunko} Гунько, А. В.    Программирование (в среде Windows) : учебное пособие / А. В. Гунько ; Новосибирский государственный технический университет. - Новосибирск : Новосибирский государственный технический университет, 2019. - 155 с. - URL: https://biblioclub.ru/index.php?page=book\&id=575417 (дата обращения: 03.05.2024) . - Режим доступа: по подписке. - Библиогр. в кн. - ISBN 978-5-7782-3890-9 : Б. ц. - Текст : электронный.
\bibitem{sravnenieNS} Черных Владимир Сергеевич, Жихарев Александр Геннадиевич, Федосеев Артемий Дмитриевич, Мартон Никита Андреевич Сравнение эффективности различных методов обучения нейронных сетей // Научный результат. Информационные технологии. 2023. №1. URL: https://cyberleninka.ru/article/n/sravnenie-effektivnosti-razlichnyh-metodov-obucheniya-neyronnyh-setey (дата обращения: 08.06.2026).
\bibitem{averch} Аверченков, В. И.    Основы математического моделирования технических систем : учебное пособие / В. И. Аверченков, В. П. Федоров, М. Л. Хейфец. - 4-е изд., стер. - Москва : Флинта, 2021. - 271 с. - URL: http://biblioclub.ru/index.php?page=book\&id=93344 (дата обращения: 16.02.2026) . - Режим доступа: по подписке. - ISBN 978-5-9765-1278-8 : Б. ц. - Текст : электронный.
\bibitem{Chistyakova} Чистякова  Т.Б., Шашихина О.Е. Интеллектуальный программный комплекс моделирования процесса планирования многоассортиментных промышленных производств // Прикладная информатика. 2022. Т. 17. № 5. С. 41–50. DOI: 10.37791/2687-0649-2022-17-5-41-50 – Текст : электронный.
\bibitem{cifradvoynik} Речкалов А.В., Артюхов А.В., Куликов Г.Г. Логико-семантическое определение цифрового двойника производственного процесса. RTJ. 2023;11(1):70–80. - URL: https://doi.org/10.32362/2500-316X-2023-11-1-70-80 (дата обращения: 16.02.2026) - Текст : электронный.
\bibitem{generativ} Виталий Романович Александров, Алексей Алексеевич Щеткин, Александр Сергеевич Бевз Генеративный искусственный интеллект в планировании производства // Измерение. Мониторинг. Управление. Контроль. 2025. №1 (51). URL: https://cyberleninka.ru/article/n/generativnyy-iskusstvennyy-intellekt-v-planirovanii-proizvodstva (дата обращения: 02.06.2026).
\bibitem{gvoz} Гвоздинский Анатолий Николаевич, Серик Екатерина Эдуардовна Исследование и разработка интеллектуальных методов управления современным предприятием // Радиоэлектроника и информатика. 2017. №3. URL: https://cyberleninka.ru/article/n/issledovanie-i-razrabotka-intellektualnyh-metodov-upravleniya-sovremennym-predpriyatiem (дата обращения: 02.06.2026).
\bibitem{rass} Рассомагин А. С. Оптимизация стратегии технического обслуживания и ремонтов с применением интеллектуальных методов // Инновации и инвестиции. 2019. №7. URL: https://cyberleninka.ru/article/n/optimizatsiya-strategii-tehnicheskogo-obsluzhivaniya-i-remontov-s-primeneniem-intellektualnyh-metodov (дата обращения: 04.06.2026).
\bibitem{yakov} Яковлева Дарья Дмитриевна Особенности системы планирования на основе интеллектуальных решений // ЭПП. 2024. №4. URL: https://cyberleninka.ru/article/n/osobennosti-sistemy-planirovaniya-na-osnove-intellektualnyh-resheniy (дата обращения: 04.06.2026).
\bibitem{alexa} Планирование производственных ресурсов на основе методов искусственного интеллекта / В.Р. Александров [и др.] // Автоматизация в промышленности. — 2024. — С. 36-39 .— Текст : непосредственный.
\bibitem{ivashen} Иващенко А. В., Никифорова Т. В. Цифровизация организационной структуры управления производственным предприятием // Известия Самарского научного центра РАН.  — 2021.  — №2. —  URL: https://cyberleninka.ru/article/n/tsifrovizatsiya-organizatsionnoy-struktury-proizvodstvennym-predpriyatiem (дата обращения: 07.06.2026).
\bibitem{preobra} Преображенский А. П., Линкина А. В. Решение задачи оптимизации ресурсоэффективности предприятия с применением технологий бережливого производства // Вестник ВГУИТ. 2021. №4 (90). URL: https://cyberleninka.ru/article/n/reshenie-zadachi-optimizatsii-resursoeffektivnosti-predpriyatiya-s-primeneniem-tehnologiy-berezhlivogo-proizvodstva (дата обращения: 05.06.2026).
\bibitem{shevcov} Шевцов Е.С., Шамин Р.В. Логическая интеграция информационных систем на основе экспертных систем. Russian Technological Journal. 2025;13(2):27-35. https://doi.org/10.32362/2500-316X-2025-13-2-27-35. EDN: SRKXBR – Текст : электронный.
\bibitem{babich} Бабич, Татьяна Николаевна. Оперативно-производственное планирование : учебное пособие / Т. Н. Бабич, Ю. В. Вертакова. - Москва : РИОР : ИНФРА-М, 2017. - 260 с. - (Высшее образование). - ISBN 978-5-369-01616-9 (РИОР). - ISBN 978-5-16-012455-1 (ИНФРА-М) - Текст : непосредственный.
\end{thebibliography}
